\begin{textsample}{POS Dim 7 – global\_south\_2025\_04 – Score 2.00 – global\_south\_2025\_04/123392106.txt}  \label{ex:f7_pos_142}
Main navigation Breadcrumb Hamas launches legal challenge against UK terror designation Apr 9, 2025 The Palestinian group Hamas has launched a legal challenge against the United Kingdom ' s 2021 decision to designate it as a terrorist organisation, submitting a detailed 106-page \textbf{application} to the Home Secretary calling for the proscription to be overturned.
The challenge is spearheaded by Dr Mousa Abu Marzouk, Head of International Relations and the Legal Office of Hamas ' Political Bureau.
His witness statement, submitted through legal representatives at Riverway Law, One Pump Court Chambers, and Garden Court Chambers, forms the basis of the \textbf{application} and includes contributions from 20 experts, including Professor John Dugard, a former judge at the International Court of Justice."Hamas is not a terrorist group,"Abu Marzouk stated in his declaration."It is a Palestinian Islamic liberation and resistance movement whose goal is to liberate Palestine and confront the Zionist project."He added,"We also look outwards to draw inspiration from the @ @ @ @ @ @ @ @ @ @ resisted colonialism, occupation and imperialism in the name of justice, dignity, and human equality."The appeal, submitted on Wednesday, argues that the proscription of Hamas not only undermines the possibility of political \textbf{resolution} in Palestine but also criminalises ordinary Palestinians and those who speak in solidarity with them.
The lawyers contend that the designation has had the effect of stifling debate and political dialogue in the UK and abroad, while preventing Hamas from engaging in conflict \textbf{resolution} efforts."The British government ' s decision to proscribe Hamas is an unjust one that is symptomatic of its unwavering support for Zionism, apartheid, occupation and ethnic cleansing in Palestine for over a century,"Marzouk added."Hamas does not and never has posed a threat to Britain, despite the latter ' s ongoing complicity in the genocide of our people.
It is perhaps out of colonial guilt that Britain fears that one day, those it oppresses will strike back against the sponsors of the Zionist entity.
Britain should have no such fear @ @ @ @ @ @ @ @ @ @ ( ANC ) and the Irish Republican Army ( IRA ), the legal team argues that proscription"undermines the possibility of a peaceful settlement.""Transition to a political process is hindered by the terrorism label, as talking with terrorists is a taboo,"Hamas said.
The UK initially banned Hamas ' s armed wing, the Qassam Brigades, in 2001.
However, in 2021 the government moved to proscribe Hamas in its entirety, claiming that distinguishing between its political and military wings was"artificial."The Home Office at the time stated that Hamas was a"complex but single terrorist organisation."Hamas does not deny that some of its actions fall under the broad definition of terrorism outlined in the UK ' s Terrorism Act 2000.
However, the group ' s legal representatives argue"that the definition also covers all groups and organisations around the world that use violence to achieve political objectives, including the Israeli armed forces, the Ukrainian Army and indeed the British armed forces."@ @ @ @ @ @ @ @ @ @, in some ways, overwhelmingly obvious,"said barrister Franck Magennis."The context is that Israel has seemingly become a pariah state and its ideology of Zionism has become toxic...
The price of associating yourself with the Israeli government is becoming far too high even for its staunchest allies."Magennis added that the Home Secretary"has an extremely broad discretion"and that there is"every reason to believe that she will find the arguments persuasive.""There ' s every reason to believe that she will find the arguments persuasive and will grant the \textbf{application} accordingly."The legal submission also criticised the chilling effect the designation has had on public debate in the UK."Rather than allow freedom of speech, police have embarked on a campaign of political intimidation and persecution of journalists, academics, peace activists and students over their perceived support for Hamas,"the lawyers argued."There is an urgent need for honest, intelligent, and nuanced conversations about the situation in @ @ @ @ @ @ @ @ @ @ of your opinion on Hamas, a policy which has the effect of stifling discussion is unhelpful and acts as a substantial hurdle to reaching a long-term political settlement."Abu Marzouk also commented on the 7 October 2023 Al-Aqsa Flood operation, stating that the objective had been a"military manoeuvre targeting the Gaza Division of Israel ' s Southern Command,"and that orders were issued to target soldiers -- not civilians."We take justice and accountability very seriously,"he wrote."We remain, as always, prepared to cooperate with any international investigations and inquiries into the operation, even if ' Israel ' refuses to do so."The UK Home Secretary now has 90 days to respond to the petition.
Should the \textbf{application} be rejected, Hamas will have the right to appeal the decision before a tribunal, which could still result in the reversal of the 2021 designation.
We need your support.
Every contribution counts.
Sri Lanka is one of the most dangerous places in the @ @ @ @ @ @ @ @ @ @ at threat, with at least 41 media workers known to have been killed by the Sri Lankan state or its paramilitaries during and after the armed conflict.
Despite the risks, our team on the ground remain committed to providing detailed and accurate reporting of developments in the Tamil homeland, across the island and around the world, as well as providing expert analysis and insight from the Tamil point of view

% matched lemmas: application, resolution
\end{textsample}
